\chapter{1966:Robert S. Mulliken}

\label{chap:chemistry1966}

The Nobel Prize in Chemistry for the year 1966 was awarded to Robert S. Mulliken.

\textbf{Motivation:}

for his fundamental work concerning chemical bonds and the electronic structure of molecules by the molecular orbital method

\section{Robert S. Mulliken}

\subsection*{Early Life and Education}

Robert S. Mulliken was born on 1896-06-07 in Newburyport, Massachusetts, USA.

\subsection*{Career and Major Research}

Mulliken studied at the Massachusetts Institute of Technology and received his doctorate from the University of Chicago in 1921. He spent most of his career at the University of Chicago, where he developed the molecular orbital method for describing chemical bonds and the electronic structure of molecules. He also studied molecular spectra and introduced an electronegativity scale based on ionization energies and electron affinities.

\subsection*{Nobel Prize-winning Work}

The work for which Robert S. Mulliken was awarded the Nobel Prize in Chemistry in 1966 was described by the Nobel Committee as: "for his fundamental work concerning chemical bonds and the electronic structure of molecules by the molecular orbital method".

Mulliken assigned electrons to molecular orbitals that extend over an entire molecule, providing a quantum-mechanical description of bonding that complemented valence bond theory. His approach explained the spectra, ionization energies, and bond properties of many molecules.

\subsection*{Significance and Impact}

The molecular orbital method became one of the two principal frameworks of quantum chemistry and is the basis of most modern computational chemistry. Mulliken's orbital diagrams and notation are still used in chemistry education.

\subsection*{Later Career and Honors}

Mulliken remained at the University of Chicago after his retirement in 1961 and held a research professorship at Florida State University from 1985. He died on 1986-10-31 in Arlington, Virginia, USA.

\subsection*{Legacy}

Mulliken's molecular orbital theory is central to modern chemistry and underpins much of computational and physical chemistry.

\subsection*{Selected Publications}

"The Assignment of Quantum Numbers for Electrons in Molecules" (Physical Review, 1928)

"Spectroscopy, Molecular Orbitals, and Chemical Bonding" (Science, 1967)

\clearpage

\section*{Chapter Summary}

The 1966 prize recognized Robert S. Mulliken for the development of the molecular orbital method. His theory of chemical bonding and molecular electronic structure remains fundamental to chemistry.

